% ============================================================ Capítulo: proceso de datos con gráficas
\chapter{Proceso de la información explicado con gráficas}\label{cap:procesodatos}

\section{Del sitio oficial a la base de conocimiento}
La información que alimenta la base de conocimiento semántica y el sistema de preguntas y respuestas proviene
del web scraping de los sitios oficiales \url{https://autopartesfanalca.com.co/} (unidad de autopartes: inicio,
nosotros, productos, clientes y contacto) y \url{https://www.fanalca.com/} (sitio corporativo). Un único comando
ejecuta el proceso completo y deja evidencia cuantitativa de cada etapa:
\begin{lstlisting}[style=ps]
cd C:\fabrica-nacional-autopartes
uv run fabrica pipeline --max-paginas 60   # rastreo → Bronze → Silver/Gold → validación → gráficas
uv run fabrica stats                       # comprobar origen "scraping" en los fragmentos
\end{lstlisting}

\begin{figure}[H]\centering
\begin{tikzpicture}[node distance=0.45cm and 0.42cm, every node/.style={font=\sffamily\scriptsize}]
  \node[externo=2.2cm] (web) {Sitios oficiales\\autopartesfanalca.com.co\\fanalca.com};
  \node[caja=1.9cm, right=of web] (r) {\textbf{1. Rastreo}\\robots.txt, sitemap, pausa $\geq$\,2\,s};
  \node[almacen, right=of r, fill=bronce!20, minimum width=1.9cm] (b) {\textbf{2. Bronze}\\HTML + manifiesto};
  \node[caja=1.9cm, right=of b, fill=plata!15] (s) {\textbf{3--4. Silver}\\extracción, limpieza, dedupe};
  \node[almacen, right=of s, fill=oro!20, minimum width=1.9cm] (g) {\textbf{5--6. Gold}\\fragmentos + métricas};
  \node[clave=2.4cm, below=1.0cm of g] (kb) {Base de conocimiento\\BM25 + LangChain};
  \node[caja=2.4cm, left=of kb] (graf) {\texttt{docs/graficas}\\seis PNG por etapa};
  \node[caja=2.6cm, below=1.0cm of r] (qa) {Resumen · FAQ · Q\&A\\API · MCP · Simulador};
  \draw[flecha] (web) -- (r);
  \draw[flecha] (r) -- (b);
  \draw[flecha] (b) -- (s);
  \draw[flecha] (s) -- (g);
  \draw[flecha] (g) -- (kb);
  \draw[flechad] (g.south west) -- node[etq, sloped, above]{métricas} (graf.north east);
  \draw[flecha] (kb.south) |- ++(0,-0.5) -| (qa.south);
\end{tikzpicture}
\caption{Flujo del proceso de datos y artefactos de cada etapa}\label{fig:flujo-proceso}
\end{figure}

\section{Dónde y cómo se ejecuta el scraping}
El scraping es un componente del sistema y se ejecuta en el computador con Windows~11, desde la consola y en la
carpeta del proyecto. No se ejecuta desde la interfaz web ni desde las pruebas automatizadas (estas usan un sitio
simulado con \texttt{responses}).
\begin{table}[H]\centering\small
\caption{Opciones para ejecutar el proceso real}\label{tab:donde-scraping}
\begin{tabularx}{\linewidth}{@{}P{3.2cm} X P{3.4cm}@{}}\toprule
Consola & Procedimiento & Control\\\midrule
PowerShell & \cmd{cd C:\textbackslash fabrica-nacional-autopartes}; \cmd{uv sync}; \cmd{uv run fabrica pipeline} & Ejecución directa por la persona responsable.\\
Claude Code (principal) & Solicitar al subagente \texttt{ingeniero-datos} (Sonnet) la fase F09; el agente propone \cmd{uv run fabrica pipeline} & Permiso \texttt{ask}: requiere aprobación porque accede a la red.\\
OpenCode (alternativa) & Mismo comando desde la sesión de OpenCode & Permiso \texttt{ask} en \texttt{opencode.json}.\\\bottomrule
\end{tabularx}
\end{table}
Tras la ejecución se versionan \ruta{data/bronze/manifiesto.jsonl}, \ruta{data/bronze/resumen_rastreo.json},
\ruta{data/gold/} y \ruta{docs/graficas/}; estos archivos son la evidencia del proceso para la presentación.

\begin{nota}[Fuente de las gráficas incluidas]
El entorno de construcción no tuvo acceso de red a los dominios de Fanalca. Las gráficas incluidas se generaron
con la base de demostración (semilla y lectura web de autopartesfanalca.com.co, con URL por fragmento) y su nota
inferior lo declara; las etapas de rastreo y Bronze muestran el aviso ``sin datos de scraping''. Al ejecutar
\cmd{uv run fabrica pipeline} en el computador local, las seis gráficas se regeneran con los datos reales.
La generación desde un rastreo se verifica en la prueba \texttt{test\_metricas\_y\_graficas\_desde\_scraping},
que recorre un sitio simulado y produce las seis gráficas.
\end{nota}

\section{Sitios con protección antibots}\label{sec:antibots}
En la primera ejecución real (4 de noviembre de 2026, Windows~11), \url{https://autopartesfanalca.com.co/}
respondió a todas las peticiones con \texttt{HTTP 202 Accepted} y el encabezado \texttt{SG-Captcha: challenge}
de SiteGround Security, mientras que \texttt{fanalca.com} entregó 60 páginas.
\begin{lstlisting}[style=txt]
$ curl.exe -sI https://autopartesfanalca.com.co/
HTTP/1.1 202 Accepted
Server: nginx
Content-Type: text/html
Content-Length: 167
SG-Captcha: challenge
X-Robots-Tag: noindex
\end{lstlisting}
El desafío no se elude (AGENTS.md, ADR-015). El rastreador lo detecta, registra el host como protegido y no lo
vuelve a solicitar; el pipeline construye la base con la fuente \texttt{mixta}: todo lo capturado por scraping
más la lectura web documentada de autopartesfanalca.com.co (\texttt{origen: lectura\_web}, con URL y fecha), solo
para los dominios sin capturas. Si la protección se retira, la siguiente ejecución rastrea el sitio y el
complemento deja de usarse sin cambios de código.

\section{Ajustes tras la primera ejecución real}
\begin{itemize}
  \item Exclusión de plantillas y páginas sin contenido (\texttt{?elementskit\_template=}, \texttt{/firma/},
        \texttt{/sitio-en-construccion/}) y de los avisos de privacidad extensos, que sumaban 12 de 90 fragmentos.
        Las exclusiones vigentes se aplican también a capturas Bronze anteriores.
  \item Clasificación de secciones por segmento exacto de la ruta: las entradas del blog se agrupan en
        \texttt{noticias} y las rutas \texttt{/negocio/} se unen a \texttt{negocios}.
  \item Opción \texttt{--desde-cero} para que el resumen y las gráficas correspondan a una ejecución completa.
  \item Segunda ejecución real (75 páginas): exclusión por expresiones regulares de la paginación de noticias
        (\texttt{/noticias/2/}) y de los filtros de vacantes por área y ciudad; el informe de sostenibilidad pasa a
        la sección \texttt{sostenibilidad} y las rutas en inglés, \texttt{/negocio/} y \texttt{holda-autos} se unifican
        en una sección por unidad de negocio. Base resultante: 201 fragmentos de scraping y 14 de lectura web.
  \item \texttt{-SoloReconstruir} reconstruye Silver, Gold y las gráficas desde el Bronze existente, sin red.
  \item Datos de contacto del pie de página: la extracción del cuerpo descarta el pie y la limpieza elimina las
        líneas repetidas, por lo que la dirección de la sede (publicada solo en el pie de fanalca.com) no llegaba a
        la base. Ahora las líneas con dirección, teléfono o correo presentes en el pie de al menos dos páginas se
        guardan una sola vez en el documento \texttt{contacto-sitio} de cada dominio.
  \item Recuperación: la pregunta «¿Dónde queda la sede principal?» no comparte vocabulario con
        «Calle 13 \# 31A - 80»; BM25 incorpora una expansión de consultas de primer contacto (ubicación y
        contacto) con peso 0,5, y los documentos de contacto se ubican primero para que el modo completo con
        recorte por presupuesto los conserve.
  \item Escritura de archivos con fin de línea LF en Windows.
\end{itemize}

\section{Métricas registradas}
El orquestador escribe \ruta{data/gold/metricas_etapas.json} con: capturas, bytes y estados HTTP de Bronze; el
resumen del rastreo (URL descubiertas en sitemaps, visitadas, capturadas, descartadas y duración); por página,
bytes de HTML, caracteres extraídos, caracteres sin plantilla, caracteres limpios, método de extracción y si se
descartó; documentos y duplicados de Silver; y fragmentos, tokens por fragmento, fragmentos por sección y por
origen en Gold. \cmd{uv run fabrica graficas} transforma ese archivo en seis figuras con una paleta sobria (azul, naranja y
aqua sobre fondo claro) y la nota de procedencia.

\section{Etapas}
\newcommand{\etapa}[3]{%
\begin{figure}[H]\centering
\IfFileExists{../graficas/#1.png}{\includegraphics[width=0.92\linewidth]{../graficas/#1.png}}{\fbox{Gráfica no generada}}
\caption{#2}\label{fig:#1}
\end{figure}
#3}

\etapa{01_rastreo}{Etapa 1: rastreo responsable}{El rastreador lee \texttt{robots.txt} de cada host, descubre
los sitemaps, respeta \texttt{Crawl-delay} y la pausa mínima, filtra dominios y rutas excluidas
(\texttt{/vacantes}, \texttt{/wp-content/uploads}) y registra cuántas URL se descubrieron, visitaron, capturaron y
descartaron.}

\etapa{02_bronze_paginas}{Etapa 2: HTML capturado por página en la capa Bronze}{Bronze conserva el HTML original
sin modificaciones, con URL final, fecha, estado HTTP y SHA-256; es la evidencia trazable del scraping.}

\etapa{03_reduccion}{Etapa 3: reducción del HTML al conocimiento}{La reducción muestra cuánto del HTML es
contenido útil: bytes de HTML, texto extraído (trafilatura con respaldo BeautifulSoup), texto sin plantilla
(menús, pies y avisos repetidos) y texto limpio y anonimizado.}

\etapa{04_limpieza_paginas}{Etapa 4: limpieza por página en la capa Silver}{Cada página se compara antes y después
de la limpieza; se descartan duplicados exactos y casi duplicados (Jaccard $\geq$ 0,9) y se anonimizan correos no
corporativos (Ley 1581 de 2012).}

\etapa{05_fragmentos_seccion}{Etapa 5: fragmentos por sección en la capa Gold}{Los documentos se fragmentan por
encabezados y tamaño con identificadores citables (\texttt{autopartes-productos\#000}); cada fragmento conserva
sección, URL y origen. Estos fragmentos son el insumo de la base semántica y del sistema Q\&A.}

\etapa{06_tokens_fragmento}{Etapa 6: tamaño de los fragmentos frente al límite de contexto}{El histograma
comprueba que ningún fragmento supera 600 tokens aproximados, límite que mantiene el contexto completo dentro de
la ventana configurada para \texttt{llama3.2:3b}.}

\section{Exploración sin notebooks}
La exploración del corpus se realiza con \cmd{uv run fabrica stats}, las gráficas por etapa y las pestañas de
Streamlit \emph{Proceso de datos} y \emph{Base de conocimiento}; el proyecto no utiliza notebooks, JupyterLab ni
Anaconda (ADR-014).
